\section{Conclusion}
\label{sec:conclusion}

QDCVR turns ``the retriever found it'' into ``the system verified it''. Its
design position is that the governed object of a knowledge-base management
system is content: content decides where a document belongs and whether a
candidate may be cited, and the absence of content that answers a question is
itself an outcome the system must report. That position is implemented as
contracts rather than conventions (RQ1): parse gates that reject visibly,
placement rules that read the document, a three-source consistency audit over
disk, tree registry and library metadata, tool whitelists enforced server-side,
and verdicts that fail closed.

The five research questions are answered by measurement rather than
assertion. Acquisition and content routing hold at the scale we tested, while
21 of the 100 submitted papers never became retrievable and the gap between the
submitted, recorded and retrievable counts is reported as an observability gap
rather than explained away (RQ2). Staged, multi-channel retrieval with content
adjudication reaches documents that one-shot vector top-$k$ misses, and the
benefit comes from staged recall plus content verification, not from the fusion
of the two retrieval stages (RQ3). Verification costs well under a second and
buys per-citation provenance, but as an ordering it inherits a same-domain
blind spot that we quantify and partially repair; the failure tracks confuser
density rather than the checkpoint, domain, language or scale (RQ4). Honest
abstention lives in the reading and answering layer, never in the threshold;
end-to-end latency belongs to model turns; and engine choice moves cost by two
orders of magnitude at the same quality tier (RQ5).

The roadmap follows from the same ledger. Verifier calibration against
per-library null distributions, trained evidence-verification models, and the
verdict-as-filter role that the tool layer already supports are the immediate
next steps; larger-corpus evaluation with an author-independent question set,
consumer-side evidence that not-found reports serve their readers, and lane
auto-selection from query shape are the open empirical ones. The direction is
towards retrieval systems whose trustworthiness is an inspectable property of
the pipeline rather than a hope about the model.
